\chapter{Glossaire}
\begin{description}
  \item[ADSR] attaque, déclin, maintien, relâchement : les quatre phases de l'enveloppe d'un son. L'attaque est la
        montée au début de la note, le déclin la descente vers le niveau de maintien, le relâchement l'extinction
        après la fin de la note.
  \item[Algorithme] sur une puce FM, la façon dont les opérateurs d'une voix sont branchés entre eux.
  \item[Arpège] accord joué en faisant défiler très vite ses notes sur une seule voix, pour donner l'illusion de
        plusieurs voix.
  \item[Cent] centième de demi-ton.
  \item[Coupure] fréquence à partir de laquelle un filtre commence à atténuer le son.
  \item[DDA] mode du HuC6280 où le processeur écrit directement les échantillons dans la voix.
  \item[Déclenchement] calage de la fenêtre de l'oscilloscope sur la période du son, pour que la forme reste
        immobile.
  \item[Driver] programme qui joue la musique sur la machine d'origine en écrivant dans les registres de la puce.
  \item[DSP] processeur de traitement du signal, utilisé pour l'écho et la réverbération.
  \item[Enveloppe] évolution du volume (ou d'un autre paramètre) d'une note dans le temps.
  \item[FIR] filtre à réponse impulsionnelle finie ; sur la Super Nintendo, il colore les répétitions de l'écho.
  \item[FM] synthèse par modulation de fréquence : un oscillateur fait varier la fréquence d'un autre, ce qui crée
        des harmoniques.
  \item[Modulation en anneau] multiplication de deux signaux, qui donne des sons métalliques ou de cloche.
  \item[Opérateur] oscillateur sinusoïdal d'une voix FM, avec sa propre enveloppe.
  \item[PAL-N] norme vidéo argentine : horloge proche du NTSC, 50 images par seconde.
  \item[Porteuse] opérateur FM que l'on entend ; les modulateurs ne s'entendent qu'à travers elle.
  \item[Rapport cyclique] proportion du temps où une onde carrée est en position haute ; il change le timbre.
  \item[Registre] case mémoire de la puce dans laquelle le driver écrit un réglage.
  \item[Résonance] accentuation des fréquences autour de la coupure d'un filtre.
  \item[Rip] fichier extrait d'un jeu, qui contient le code et les données de sa musique.
  \item[Synchronisation] remise à zéro d'un oscillateur à chaque cycle d'un autre, qui donne des sons durs et
        riches.
  \item[Table d'ondes] forme d'onde mémorisée, rejouée en boucle.
  \item[Tracker] logiciel de composition par motifs, et format des musiques qu'il produit.
  \item[Voie] ligne de l'oscilloscope qui montre une voix.
\end{description}
